\section{Worum es geht}

\textbf{[SETZUNG]} FFGFT beschreibt die fundamentale Physik auf einem kompakten
4-Torus $T^4$ mit dem einen Parameter $\xi=4/30000$. Die beobachtbare Welt ist
$\mathbb{R}^3\times\mathbb{R}_t$. Dieses Dokument bearbeitet die Frage, wie beides
zusammenhängt -- und zeigt, dass die Kette $T^4\to T^0$ aus \emph{zwei qualitativ
verschiedenen} Operationen besteht, die nicht verwechselt werden dürfen: einer
\emph{Dualitäts-Dekompaktifizierung} (fast reversibel, keine Modenverluste) und
einer \emph{geometrischen Projektion} (verlustbehaftet, nicht reversibel). Dazu
tritt eine dritte Operation, die in die andere Richtung geht: die
\emph{repräsentationelle Übersetzung} auf den Hilbertraum (bijektiv, verlustfrei).

Das Standardmodell ist darin als \emph{dekompaktifizierte Projektion} eingeordnet:
seine Konstanten ($\alpha$, $v$, Yukawa) sind Rekursions-Ausgänge aus $\xi$; wer
sie als Eingaben nimmt, arbeitet im ausgerollten Bild, ohne das kompakte Ganze zu
benennen. Das ist keine Widerlegung des Standardmodells -- es rechnet korrekt in
seinem Bereich; es ist eine Verortung.

\section{Die Spannung in der Fragestellung}

\textbf{[B]} Drei Beobachtungen zum Übergang $T^4\to\mathbb{R}^3\times\mathbb{R}_t$
stehen scheinbar im Widerspruch:

\begin{itemize}[leftmargin=2em,itemsep=3pt,topsep=3pt]
  \item \textbf{Verlustbild.} Jede Projektion $D\to D-1$ ist nicht injektiv,
    verliert Information. Die Kette ist nicht reversibel.
  \item \textbf{Kompaktheitsbild.} $T^4$ ist vollständig kompakt; die
    beobachtbare Raumzeit ist dekompaktifiziert.
  \item \textbf{Aufwärtsbild.} Es gibt bereits eine \emph{bijektive} Übersetzung
    von $T^4$ in einen höherdimensionalen Raum -- den Hilbertraum (A070) --, die
    \emph{verlustfrei} ist.
\end{itemize}

\textbf{[B]} Diese drei Bilder lassen sich nicht naiv übereinanderlegen: Wäre der
erste Schritt reiner Verlust, könnte daraus keine dekompaktifizierte Zeitrichtung
entstehen (Vergessen erzeugt keine neue Dimension); und es bliebe rätselhaft, wie
eine Aufwärtsbewegung verlustfrei sein kann, während die Abwärtsbewegung
Information vernichtet. Die Auflösung: Es gibt nicht eine, sondern
\emph{drei} dimensionsändernde Operationen von \emph{zwei Arten}.

\section{Die Struktur von $T^4$}

\textbf{[B]} Die fundamentale Ebene der FFGFT ist der kompakte 4-Torus
\[
  T^4 = S^1_x\times S^1_y\times S^1_z\times S^1_m,
\]
mit drei \emph{räumlichen} Kreisen und einem \emph{Masse-Kreis} $S^1_m$. Die
vierte Richtung $S^1_m$ ist nicht eine Massedimension \emph{neben} einer
separaten Zeit, sondern die \emph{Ersatz-Zeit-Richtung} selbst. Masse und Zeit
sind zwei \emph{Lesarten} ein und derselben Dimension; die Fundamentalrelation
$\tilde T\cdot m=1$ ist das Übersetzungswörterbuch:
\[
  \underbrace{S^1_m\;\text{(kompakt)}}_{\text{Masse-Lesart, diskretes Spektrum}}
  \;\;\xleftrightarrow{\;\;\tilde T\cdot m=1\;\;}\;\;
  \underbrace{\mathbb{R}_t\;\text{(offen)}}_{\text{Zeit-Lesart, kontinuierlicher Fluss}}.
\]

\textbf{[K]} \textbf{Kaluza-Klein-Lesart.} Ein kompakter Kreis vom Radius $R$
trägt einen diskreten Massenturm $m_n\sim n/R$. Die vier Richtungen tragen zwei
Radienklassen: die drei Raumkreise ($R_{x,y,z}\sim R_\text{groß}$) erzeugen einen
quasi-kontinuierlichen Moden-Turm mit $\sim0$ Lücke -- das wird zum Raum; der
Massekreis ($R_m\sim\xi\,\ell_P$) hat eine riesige Lücke und ein diskretes,
schweres Spektrum -- das wird zum Massenspektrum der Teilchen. Die
Fundamentalrelation $\tilde T\cdot m=1$ ist damit strukturell die
Kaluza-Klein-Relation für die Massedimension.

\section{Typ I -- Dualitäts-Dekompaktifizierung: Masse wird Zeit}

\textbf{[B]} Die erste, qualitativ ausgezeichnete Operation ist das \emph{Abrollen}
des Massekreises über seine universelle Überlagerung:
\[
  S^1_m\;\xrightarrow{\;\text{Überlagerung}\;}\;\mathbb{R}_t,
  \qquad p:\mathbb{R}_t\to S^1_m,\;\;t\mapsto e^{2\pi i\,t/R_m},
\]
gesteuert durch $\tilde T\cdot m=1$. Diese Operation ist \emph{einbettend}, keine
Projektion: sie erzeugt eine neue Richtung ($\mathbb{R}_t$), statt eine
wegzulassen. Periodische Strukturen auf $S^1_m$ heben sich eindeutig in
$\mathbb{R}_t$; kein Modeninhalt geht verloren. Der \enquote{Preis} der
Dekompaktifizierung ist kein Informationsverlust, sondern das Auftreten der
Zeitrichtung selbst.

\textbf{[B]} Was nicht direkt sichtbar bleibt: die \emph{Kompaktheit}. In der
Masse-Lesart ist die Topologie manifest -- das Spektrum ist diskret, die
Identifikation $t\sim t+R_m$ direkt ausgedrückt. In der Zeit-Lesart ist dieselbe
Kompaktheit nur \emph{latent}: aus kontinuierlicher Zeit allein ist die
topologische Identifikation nicht eindeutig rekonstruierbar (die Windung bleibt
mehrdeutig). Genau dort setzt die Rekursion $\xi^n$ an -- sie \emph{ist} die
Spektralsignatur, an der die verborgene Kompaktheit erkennbar bleibt.

\textbf{[S]} \textbf{Was die ausgerollte Zeit stetig und gleichförmig macht.}
Drei Fragen, drei Antworten: (1) \emph{Warum ist $\mathbb{R}_t$ stetig?} Weil das
Abrollen die universelle Überlagerung $p:\mathbb{R}\to S^1$ ist -- ein lokaler
Homöomorphismus; die Überlagerung eines Kreises \emph{ist} die stetige Gerade.
(2) \emph{Warum fließt die Zeit gleichförmig?} Wegen der Homogenität von $S^1_m$
(translationsinvariant, keine ausgezeichneter Punkt) und des einen festen Eigenwerts
$m=1/\tilde T$. (3) \emph{Warum wird überhaupt ausgerollt?} Nicht die Geometrie
entscheidet das -- im kompakten $T^4$ sind alle vier Kreise gleichberechtigt, es
gibt keine Zeit. Der \emph{Messvorgang} (der eingebettete Beobachter) wählt einen
Kreis als Zeit und rollt ihn auf. Die Stetigkeit ist die vergessene Windung: behielte
man sie, hätte man den diskreten Kreis zurück.

\section{Typ II -- geometrische Projektion: Raumrichtungen vergessen}

\textbf{[B]} Alle weiteren Schritte ($D3\to D2\to D1\to D0$) sind echte
geometrische Projektionen räumlicher Richtungen -- eine Projektion
$\pi:\mathbb{R}^k\to\mathbb{R}^{k-1}$ mit $\ker(\pi)=\mathbb{R}$. Sie kollabiert
eine vollständige Faser; die Information über die vergessene Richtung ist im Bild
nicht enthalten. Ohne Zusatzdaten (holographische Vollständigkeit,
Randbedingungen) nicht rekonstruierbar. Die Rekursion $\xi^n$ wirkt auf diese
Schritte \emph{nicht} -- räumliche Projektionsverluste sind echt.

\textbf{[B]} Der Schritt $D3\to D2$ ist mathematisch das holographische Prinzip;
im Horizontfall erscheint hier der Bekenstein-Hawking-Bound $S\le A/(4\ell_P^2)$.

\section{Typ III -- repräsentationelle Übersetzung: bijektiv nach oben}

\textbf{[B]} Abwärts (geometrisch) und aufwärts (repräsentationell) sind
\emph{entgegengesetzte} Informationseigenschaften:
\[
  T^4 \;\xleftrightarrow{\;\text{bijektiv}\;}\; \mathcal H = L^2(T^4)\otimes\mathbb{C}^3
\]
Die Übersetzung (A070) ist keine Approximation, sondern eine bijektive
Strukturaussage: jede QM-Rechnung ist in FFGFT-Sprache abbildbar und umgekehrt.
Es werden \emph{keine} physikalischen Extra-Dimensionen hinzugefügt -- derselbe
4D-Inhalt wird in einer höherdimensionalen \emph{Sprache} ausgedrückt, in der er
ohne Reduktion Platz hat. Die Verlustfreiheit folgt gerade aus dem
repräsentationellen Charakter.

\textbf{[B]} Eine geometrische Erweiterung $T^4\to T^5\to\dots$ ist in FFGFT
ausgeschlossen: Komplexität entsteht aus \emph{fraktaler Tiefe}
($D_f=3-\xi$, Selbstähnlichkeit), nicht aus \emph{mehr} Dimensionen. Das ist der
bewusste Gegensatz zur Stringtheorie (10/11D, Calabi-Yau,
Landschaftsproblem). \enquote{Höher} heißt in FFGFT stets Typ III
(repräsentationell), nie geometrisch.

\section{Die formale Kette, Stufe für Stufe}

\textbf{[B]} Die Kette $T^4\to T^0$ nach Typ getrennt:

\begin{center}
\renewcommand{\arraystretch}{1.35}
{\small
\begin{tabular}{lllll}
\toprule
\textbf{Stufe} & \textbf{Operation} & \textbf{Typ} & \textbf{Invariant} & \textbf{$\xi^n$} \\
\midrule
$D4\to D3{+}t$ & $S^1_m\to\mathbb{R}_t$ (abrollen) & I & $\tilde Tm=1$ & Spektralsignatur \\
$D3\to D2$ & Raum vergessen & II & Fläche $A$ & --- \\
$D2\to D1$ & Raum vergessen & II & Länge $L$ & --- \\
$D1\to D0$ & Raum vergessen & II & $n\in\mathbb{Z}$ & --- \\
\bottomrule
\end{tabular}
}
\renewcommand{\arraystretch}{1.0}
\end{center}

\[
\underbrace{T^4\;\longrightarrow\;\mathbb{R}^3\times\mathbb{R}_t}_{\text{Typ I: einbettend, fast reversibel}}
\;\longrightarrow\;
\underbrace{\mathbb{R}^2\;\longrightarrow\;\mathbb{R}\;\longrightarrow\;\{\mathrm{pt}\}}_{\text{Typ II: projizierend, nicht reversibel}}
\]

\textbf{[B]} Stufe~1 bündelt zwei informationserhaltende Operationen: (Ia) der
Massekreis rollt über die Dualität $\tilde Tm=1$ in $\mathbb{R}_t$ ab; (Ib) die
drei Raumkreise dekompaktifizieren durch den \emph{Grenzprozess} großer Radien
($R_{x,y,z}\to\infty$) zu $\mathbb{R}^3$ -- ein Limes, kein
Dualitätsschritt, aber ebenfalls einbettend. Beide sind einbettend; der
Mechanismus ist verschieden.

\section{Drei Operationen im Überblick}

\textbf{[B]}
\begin{center}
\renewcommand{\arraystretch}{1.35}
{\small
\begin{tabular}{llll}
\toprule
\textbf{Typ} & \textbf{Operation} & \textbf{Charakter} & \textbf{Information} \\
\midrule
I & Masse\,$\leftrightarrow$\,Zeit ($\tilde Tm=1$) & geometrisch & fast reversibel \\
II & räumliche Projektion (runter) & geometrisch & verlustbehaftet \\
III & $T^4\leftrightarrow\mathcal H$/Hilbertraum (hoch) & repräsentationell & \textbf{bijektiv} \\
\bottomrule
\end{tabular}
}
\renewcommand{\arraystretch}{1.0}
\end{center}

Die Bijektivität von Typ~III ist die spiegelbildliche Aussage zur
Nicht-Reversibilität von Typ~II: geometrisch verliert man beim Reduzieren,
repräsentationell verliert man beim Übersetzen nichts. Bild~(A) und Bild~(B) der
Eingangsspannung sind beide korrekt -- sie betreffen verschiedene Operationen.

\section{Die Rolle der Rekursion $\xi^n$}

\textbf{[K]} Die Rekursion $\xi^n$ ist \emph{keine} universelle Kompensation für
Projektionsverluste -- ihre Funktion ist eng begrenzt. Bei Typ~I erzeugt $\xi^n$
das diskrete Massenspektrum (Teilchenleiter) -- das ist der kompakte Modeninhalt
von $S^1_m$, dual als Zeit-Frequenz-Struktur: $\xi^n$ \enquote{füllt} nichts auf,
sondern ist die Spektralsignatur, an der die verborgene Kompaktheit erkennbar
bleibt. Bei Typ~II wirkt $\xi^n$ nicht -- räumliche Projektionsverluste sind echt
und werden von der Rekursion nicht berührt. Eine Darstellung, die der Rekursion
eine über alle vier Stufen monoton abnehmende \enquote{Kompensationskraft}
zuschreibt, ist irreführend.

\section{Reversibilität, präzise pro Typ}

\textbf{[B]} Der Typ-I-Schritt ist als Dualität (Fourier/KK) im Modeninhalt
umkehrbar; die Windung bleibt mehrdeutig -- \enquote{fast reversibel}. Die
räumliche Teilkette $\mathbb{R}^3\to\{\mathrm{pt}\}$ ist nicht invertierbar: es
gibt keine stetige Abbildung $\rho:\{\mathrm{pt}\}\to\mathbb{R}^3$ mit
$\rho\circ\Pi=\mathrm{id}$ -- jedes $\pi_k$ hat $\ker=\mathbb{R}\neq\{0\}$, die
kollabierte Faserstruktur ist im Bild nicht mehr vorhanden.

\section{Warum genau vier Kreise}

\textbf{[B]} Keine Theorie leitet die Dimensionszahl der Raumzeit her; auch FFGFT
\emph{leitet die 4 nicht vorwärts her} -- sie nimmt $3+1$ als empirische Eingabe.
Die Frage ist daher nicht \enquote{warum ist die Welt vierdimensional?}, sondern:
\emph{Gegeben $3+1$ -- welche Struktur ist die minimale, die alles trägt?}

\textbf{[K]} Naiv bräuchte man \emph{fünf} Bausteine: drei für den Raum, einen
für ein Massenspektrum, einen für die Zeitentwicklung. Die Fundamentalrelation
$\tilde T\cdot m=1$ identifiziert jedoch Masse-Lesart und Zeit-Lesart \emph{einer}
Dimension (Typ~I): der vierte Kreis $S^1_m$ trägt als KK-Turm das Massenspektrum
\emph{und}, dekompaktifiziert, die Zeit. Die Dualität lässt einen Kreis die Arbeit
von zweien tun -- die naive $5$ kollabiert auf $\mathbf{4}$. Weniger als vier
($T^3$ allein): kein Massenspektrum, keine Zeit. Mehr als vier (ein zweiter
Dual-Kreis): ein zweiter unabhängiger Massenturm -- überzählig, ohne Korrelat.
\enquote{Optimal} heißt hier \emph{minimal-suffizient}: die kleinste
Dimensionszahl, in der Raum, volles Massenspektrum und Zeit gemeinsam aus einem
geometrischen Objekt entstehen.

\textbf{[B]} Die Vier-Kreis-Struktur ist zugleich die Exponenten-Buchhaltung der
SI-Umrechnungsfaktoren: $c$ ist die Raum-Zeit-Brücke der drei Raumkreise, $\hbar$
die Masse-Zeit-Brücke des Massekreises. Natürliche Einheiten ($c=\hbar=1$)
kollabieren alle drei Dimensionen auf eine Achse und vernichten genau die
Unterscheidung, die die Exponentenstruktur trägt; $c$ und $\hbar$ dürfen erst
\emph{nach} dem Ablesen der Exponenten auf $1$ gesetzt werden.

\section{Im kompakten Bild: der Zeit-Kreis wird durch den Messvorgang gewählt}

\textbf{[S]} Im kompakten $T^4$ sind alle vier Kreise identisch: dimensionslose
zyklische Koordinaten, keine Signatur, keine Raum/Zeit-Unterscheidung. Die
Asymmetrie (eine Zeit, Signatur $(3,1)$) erscheint erst unter der
Dekompaktifizierung (Typ~I): genau ein Kreis rollt durch seinen Universalüberzug in
die nichtkompakte Zeit aus, die anderen drei bleiben kompakt (diskretes Spektrum)
oder projizieren (Raum). \emph{Welcher} Kreis zur Zeit wird, ist keine
intrinsisch-geometrische Tatsache -- die Geometrie ist symmetrisch. Zeit ist
\emph{relational}: die Auswahl ist der Messvorgang selbst. Der eingebettete
Beobachter, der von innen ausliest, erlebt eine Richtung als Dauer;
dieses Erleben konstituiert den Zeit-Kreis. Die Symmetriebrechung ist der
Messvorgang, nicht die Geometrie.

\textbf{[B]} \textbf{Drei Aussagen schließen die Buchhaltung lückenlos:} (1) Dass
es \emph{eine} Zeit gibt (nicht zwei): das ist der Signatur-Constraint -- eine
zeitartige Richtung gibt hyperbolische Feldgleichungen mit wohlgestelltem
Cauchy-Problem; zwei oder mehr geben ultrahyperbolische Gleichungen ohne
wohlgestelltes Anfangswertproblem. (2) \emph{Welcher} Kreis die Zeit ist: durch
den Messvorgang festgelegt (die erlebte Situation des eingebetteten Beobachters).
(3) Dass die Aufteilung $3+1$ ist: die Struktur des Eingebettet-Seins -- die
empirische Eingabe.

\section{Die topologische Identifikation $t\sim t+R_m$ als physikalische Observable}

\textbf{[K]} Die topologische Identifikation ist keine reine Begleitstruktur,
sondern eine physikalische Observable. In der Zeit-Lesart ist ihre Periode die
fundamentale Umlaufzeit
\[
  T_\text{rev} = \frac{R_m}{c} = \frac{\xi\,\ell_P}{c}\approx 7{,}2\times10^{-48}\,\text{s},
  \qquad \Delta\phi = 2\pi\xi\approx 8{,}4\times10^{-4}\;\text{rad/Umlauf}.
\]

\textbf{[K]} Zwei prüfbare Lesarten: (1) \textbf{Masse-Lesart (reif,
quantitativ).} Das diskrete Modenspektrum erzeugt die Teilchenmassen als
Torus-Obertöne; daraus folgen Koide, $g-2$ und drei Generationen (A100/A110/A120).
Die Kompaktheit ist zahlenmäßig eingelöst. (2) \textbf{Zeit-Lesart
(falsifizierbar).} Die Periodizität erzeugt eine CHSH-Abweichung von der
Tsirelson-Schranke der Ordnung $\xi$; für 73 Qubits $\Delta\approx5{,}4\times10^{-4}$
(auflösbar bei $\sigma\sim10^{-4}$). Der Hardware-Test IBM Heron r2 (Mai 2026)
bestätigt die Bell-Verletzung ($S=2{,}74$); der $\xi$-Effekt liegt unter dem
NISQ-Rauschen und ist noch nicht auflösbar (A160). Keine unfalsifizierbare
Begleitstruktur.

\textbf{[S]} \textbf{Skalierung: linear vs.\ logarithmisch.} Eine naive
\emph{kohärente Linearakkumulation} $r\cdot\Delta\phi$ über $r$ Umläufe ergäbe
physikalisch absurde Phasendriften ($\sim10^{44}$ rad/s) und stünde im Widerspruch
zu funktionierenden Atomuhren und tiefen Schaltkreisen. Die Labor-Vorhersage
skaliert \emph{logarithmisch},
\[
  \delta\phi\sim\frac{\xi\,\ln(r)}{D_f}\approx2\times10^{-4}
  \quad (D_f=3-\xi),
\]
nicht linear. Die lineare Form ist die Worst-Case-Schranke des RSA-Arguments
(Kohärenz über $r\approx N$ Zyklen), nicht die Laborvorhersage. Diese
Regime-Trennung erklärt, warum tiefe kohärente Schaltkreise laufen.

\section{Das Standardmodell als dekompaktifizierte Projektion}

\textbf{[K]} Jede Rechnung, die $\alpha$, $v$ oder Yukawa-Kopplungen verwendet,
führt aus FFGFT-Sicht die $K_\text{frak}$-Rekursion \emph{indirekt} mit -- denn:
\begin{itemize}[leftmargin=2em,itemsep=2pt,topsep=2pt]
  \item $\alpha=\xi\,E_0^2$ (A130), wobei $E_0$ den $K_\text{frak}$-Faktor bereits
    absorbiert enthält. Wer $\alpha$ verwendet, verwendet einen Wert, in den die
    Rekursion eingegangen ist.
  \item Die Massen sind die $\xi^p$-Leiter ($m=r\,\xi^p\,v$, A100); eine
    Yukawa-Kopplung $g=m/v$ ist damit bereits ein $\xi$-Leiter-Verhältnis. Wer
    Yukawa schreibt, schreibt ein Rekursions-Verhältnis in anderer Notation.
\end{itemize}

Das Standardmodell nimmt diese Größen als gemessene Eingaben; FFGFT \emph{leitet}
sie ab. Wer $\alpha$, $v$ oder Yukawa in kontinuierlichen Feldern auf
$\mathbb{R}$ verwendet, arbeitet strukturell im \emph{ausgerollten} Bild -- dem
Typ-I-Bereich der Dekompaktifizierung. Die kompakte Herkunft der Konstanten ist
mitgeführt, aber nicht benannt. Das ist dasselbe Muster wie eine Projektion, die
korrekt auswertet, aber das Objekt nicht benennt, dessen Ausschnitt sie ist.

\textbf{[S]} \textbf{FFGFT als verbindendes Erklärungsmodell.} FFGFT verbindet QM
und Standardmodell nicht dadurch, dass es sie ersetzt, sondern dadurch, dass es
beide als \emph{Projektionen desselben kompakten Objekts} einordnet. QM und RT
sind Projektionen auf verschiedene Skalen; das Standardmodell ist die
dekompaktifizierte Ein-Skalen-Projektion seiner Konstanten. Daraus folgt eine
Eigenschaft, die FFGFT unter bekannten Rahmen auszeichnet: jedes auf seiner Skala
korrekt rechnende Teilmodell ist im Prinzip als Projektion einbaubar -- kein
Konkurrent, sondern eine Verortung. Das funktionierende Teilmodell wird nicht
widerlegt; FFGFT gibt den Ort an, an dem es sitzt.

\textbf{[S]} \textbf{Die Einordnung ist bedingt.} Die Aussage \enquote{das
Standardmodell ist eine dekompaktifizierte Untermenge von FFGFT} ist eine
FFGFT-interne Einordnung, kein von außen erzwingbarer Beweis. Sie gilt
\emph{unter der Annahme}, dass FFGFTs Ableitungen zutreffen: \emph{wenn}
$\alpha=\xi E_0^2$ und die $\xi^p$-Massenleiter korrekt sind, \emph{dann} ist das
Standardmodell die ausgerollte Projektion. Das Standardmodell kann die Relation
weder widerlegen (es rechnet korrekt im ausgerollten Bild) noch von innen
bestätigen (es benennt kein kompaktes Objekt).

\textbf{[S]} \textbf{Was belegt ist, was Programm bleibt.} Belegt ist die
Enthaltensein-Relation für den Sektor, den FFGFT tatsächlich ableitet: $\alpha$,
die Leptonmassen, der Koide-Skalar $Q=2/3$, die Phase $\theta=2/9$ (A120/A130) --
für diese ist gezeigt, dass sie Rekursions-Ausgänge sind. Programm bleibt die
\emph{volle} Yukawa-Matrix: Quarkmassen, Mischungswinkel (CKM/PMNS), CP-Phase.
Diese Trennung ist dieselbe wie im gesamten Korpus: Kern-Ableitungen (aus $\xi$
bewiesen), Brücken (algebraisch bewiesen), Reduktionen (Plausibilitätsskizzen).

\section{Erweiterung innerhalb der Projektion -- und darüber hinaus}

\textbf{[S]} In jeder Projektion kann man \emph{verfeinern}, aber nicht über
ihren Rand \emph{hinaus}. Eine Projektion hat das kompakte Ganze bereits
weggelassen; eine Erweiterung innerhalb der Projektion kann das Weggelassene
nicht zurückholen. Das ist die Geschichte des Standardmodell-Lagrangian: tritt
eine Unvollständigkeit auf -- Neutrinos massiv, eine Anomalie, ein Sektor passt
nicht --, stückelt man an: ein weiterer Term, ein weiteres Feld. Der Lagrangian
wächst, bleibt aber innerhalb derselben Projektion. Jede solche Erweiterung muss
ihre Bausteine \emph{importieren} ($\alpha$, $v$, Massenverhältnisse -- im
Teilbereich nicht ableitbar, weil sie zum weggelassenen Ganzen gehören). Der
Import ist die Signatur der Erweiterung innerhalb der Projektion; die Ableitung
ist die Signatur des benannten Ganzen.

\textbf{[S]} \textbf{Die vor-kopernikanische Analogie.} Die
vor-kopernikanische Astronomie war nicht falsch (die Epizykel-Astronomie sagte
Planetenpositionen präzise voraus) -- sie war eine \emph{innere Sicht}: die
Perspektive eines Beobachters, der seinen Standort für das Zentrum hält. Bei
Abweichungen stückelte man an -- ein weiterer Epizykel. Das Standardmodell folgt
demselben Muster: die innere Sicht ist die Projektion, die die eigene Skala für
das Ganze nimmt; das Anstückeln von Termen ist das Anstückeln von Epizykeln;
präzise aber nicht erklärend. Der kopernikanische Schritt war kein besserer Fit
-- anfangs war er schlechter (Kopernikus behielt Kreisbahnen). Er war ein
\emph{Wechsel der Perspektive}: das System als Ganzes benannt, aus dem die
scheinbaren Bewegungen als Projektionen folgen. Genau so verhält sich FFGFT zum
Standardmodell: die scheinbare Willkür der 19~Parameter ist aus der inneren Sicht
ein Satz anzupassender Zahlen; aus der äußeren sind es Projektionseffekte einer
Rekursion ($\alpha=\xi E_0^2$, $\xi^p$-Massenleiter). Zwei Grenzen halten die
Analogie ehrlich: sie beschreibt die \emph{Struktur} der beiden Sichtweisen,
garantiert aber nicht den historischen Ausgang; und die beliebige
Anstückelbarkeit des Standardmodells ist das Merkmal der inneren Sicht -- sie kann
nie von innen scheitern, nur von außen überholt werden.

\section{Was offen bleibt}

\textbf{Erstens ist die großräumige Außenskala in diesem Dokument als
deklarierte externe Eingabe geführt, nicht aus $\xi$ hergeleitet.} Die
Zwei-Typen-Struktur hängt nur an der Radienklassen-Trennung groß/klein, nicht am Zahlenwert. Die Verbindung der großräumigen Außenskala zu $\ell_1$ bleibt offen.

\textbf{Zweitens ist die holographische Vollständigkeit bei Typ II ($D3\to D2$)
nicht ausgeführt.} Inwieweit lässt sich der Typ-II-Verlust über Randdaten exakt
(nicht nur teilweise) rekonstruieren?

\textbf{Drittens ist die volle Yukawa-Matrix Programm.} Die SM-Untermengen-Relation
ist für den Leptonsektor Kern-Ableitung, für Quarkmassen/CKM/PMNS/CP-Phase offene
Bringschuld.

\section{Prüfskripte}

\begin{itemize}[leftmargin=2em,itemsep=2pt,topsep=2pt]
  \item Typ-I-Reversibilität: Fourier-Bijektion $S^1_m\leftrightarrow\mathbb{R}_t$
    numerisch, und Windungsmehrdeutigkeit:
    \texttt{python/A\_Serie\_Skripte/a090\_projektion.py}
  \item CHSH-Abweichung: $\Delta\text{CHSH}\approx\xi/(2\pi)$ numerisch und
    logarithmische vs.\ lineare Regime-Skalierung:
    \texttt{python/A\_Serie\_Skripte/a090\_chsh.py}
\end{itemize}

\section{Ersetzt}

Dieses Dokument nimmt auf: Dok.~270 (Projektionskette $T^4\to T^0$, Typ I/II/III,
Reversibilität, Rolle $\xi^n$, warum vier Kreise, Zeit-Kreis als Messvorgang,
$t\sim t+R_m$ als Observable, lineare/logarithmische Skalierung), Dok.~298 (SM als
dekompaktifizierte Projektion, Kopernikus-Analogie, was belegt/Programm).
\emph{Nicht} übernommen aus Dok.~263 (Fraktale Holografie): der Großteil mit 31
Kosmologie-Treffern, insbesondere $H_0$ als Projektions-Artefakt, Vakuumenergie,
UIFT-Brücke -- alle gehören nicht in den A-Serien-Kern (A010). Ebenfalls nicht
übernommen aus Dok.~298: der HLV-YUKAWA-Abschnitt (IPI-Rahmen). Die
holographische Projektion $D3\to D2$ (Dok.~263/268/254) ist benannt, aber nicht
ausgeführt (offener Punkt).

\section{Verwendung von KI-Werkzeugen}

Dieses Dokument ist unter Verwendung KI-gestützter Werkzeuge entstanden. Die
Fragestellungen, die Setzungen und alle inhaltlichen Entscheidungen stammen vom
Autor; die Werkzeuge formulieren aus, rechnen nach und erstellen Prüfskripte.
Die Verantwortung für Inhalt und Fehler liegt vollständig beim Autor. Der
ausführliche Wortlaut steht in A010.
